\section{Hessian 与无约束极值}
\label{sec:02-Calculus/06-Multivariable-Calculus/05}

\subsection{梯度为零后，答案分岔了}

你找到了梯度为零的点。在单变量微积分里，下一步看二阶导数的正负——正则是极小，负则是极大。搬到多元，你直觉上也想这么干。

拿 \(f(x,y)=x^2-y^2\) 在原点试。\(\nabla f(0,0)=(0,0)\)。看沿 \(x\) 轴：\(f(x,0)=x^2\)，二阶导 \(2>0\)，像是极小。看沿 \(y\) 轴：\(f(0,y)=-y^2\)，二阶导 \(-2<0\)，像是极大。同一个点，不同坐标轴给出相互矛盾的结论。沿 \(y=x\)：\(f(t,t)=0\)，恒为零，不增不减。

原点既不是极小也不是极大——它是一个鞍点。单变量那个干净的二阶导测试，搬到多元直接分裂了。

\subsection{不是一两个方向的曲率，是所有方向的曲率}

单变量里，函数局部行为只沿数轴的正负两个方向展开。二阶导数的符号同时决定了这两个方向上的凸性——一致是极小，一致是极大，没空间容纳第三种可能。

二维以上，方向多了。沿每个方向 \(f(\mathbf a+t\mathbf u)\) 都是一根独立的单变量函数，各有各的二阶导数。如果有些方向弯曲向上、有些弯曲向下，梯度为零的点就什么都不是——这便是鞍点。

所以多元极值的二阶判定，不能只看一两个坐标方向上的曲率。你需要确认函数在该点附近的每一个方向上都弯向同一侧。这等价于检查一个二次型的正定性。

\subsection{二阶 Taylor 展开：Hessian 是二阶局部模型的心脏}

对二阶连续可微的标量函数，展开到二阶：
\[
f(\mathbf a+\mathbf h)
=f(\mathbf a)
+\nabla f(\mathbf a)^{\trans}\mathbf h
+\frac12\mathbf h^{\trans} H_f(\mathbf a)\mathbf h
+o(\|\mathbf h\|^2),
\]
其中 Hessian 矩阵为
\[
H_f(\mathbf x) = \begin{bmatrix}
\partial^2f/\partial x_i\partial x_j
\end{bmatrix}_{ij}.
\]

若混合偏导在附近连续，Clairaut 定理保证 \(f_{x_ix_j}=f_{x_jx_i}\)——Hessian 是对称矩阵。这个对称性不是碰巧，而是光滑性的一个深层后果。

梯度项 \(\nabla f^{\trans}\mathbf h\) 管局部倾斜；Hessian 二次型 \(\frac12\mathbf h^{\trans}H\mathbf h\) 管不同方向上的弯曲程度。因为是实对称矩阵，谱定理保证 Hessian 可以沿一组互相正交的特征方向对角化——二次项拆成独立的平方项之和（线性代数部分在``特征值与对角化''一章有完整推导）。

在驻点 \(\nabla f(\mathbf a)=0\) 处，一阶项消失，局部行为由 Hessian 二次型决定。若对每一个非零方向 \(\mathbf h\)，都有 \(\mathbf h^{\trans}H\mathbf h>0\)，则 Hessian 正定，所有方向弯曲向上——严格局部极小。若对所有方向都有 \(\mathbf h^{\trans}H\mathbf h<0\)，负定，严格局部极大。若有的方向正、有的方向负，不定，鞍点。

\subsection{咖啡配方：当 Hessian 半正定时会发生什么}

调咖啡粉量 \(m\) 和水 \(w>0\)，只用浓度偏差来评分：
\[
f(m,w)=\left(\frac mw-c\right)^2.
\]
所有满足 \(m=cw\) 的配方都达到同样的最低分（零）。沿着 \(m\) 和 \(w\) 等比例增加的方向——比如咖啡粉和水各加一倍——浓度不变，评分不变。这是一条平坦的谷底，不是孤立的极小点。

在谷底上的任意一点，梯度为零。Hessian 为
\[
\frac{2}{w^2}
\begin{pmatrix}1&-c\\-c&c^2\end{pmatrix}.
\]
矩阵的行列式为零——有一个零特征值，对应切向量 \((c,1)^{\trans}\)。沿这个方向，Hessian 二次型恒为零，二阶信息完全不给任何曲率判断。正交方向上的曲率为正——跨出谷地改变浓度，评分立刻上升。

这就是半正定的处境：没有方向弯曲向下，但至少有一个方向二阶信息为零。二阶检验没有结论——极小可能存在（如此例），也可能不存在。需要看更高阶项或直接分析原函数。

在模型校准中，参数补偿产生的平坦方向是同一结构：训练损失对某些参数组合不敏感，因为多个参数联合变化时输出不变。``找到梯度为零的解''不等于每个参数都被数据可靠地确定。

比较曲率时还要注意变量的单位和缩放：以克为单位和以毫升为单位，坐标刻度本就不同，Hessian 的特征值数值也随之变化。曲率的实际意义需要在同一个物理尺度下比较。

\subsection{正定性的判定：特征值还是主子式}

对 \(n\times n\) 实对称矩阵，判断正定性有两条等价道路：
\begin{itemize}
\tightlist
\item 所有特征值大于零；
\item 所有顺序主子式大于零（Sylvester 准则）。
\end{itemize}

对二元矩阵，顺序主子式就是 \(f_{xx}\) 和行列式 \(D=f_{xx}f_{yy}-f_{xy}^2\)。记结论背后的几何比分记公式更可靠：正定意味着椭圆——沿任何方向，二次型都给正值；不定意味着双曲面——有的方向给正、有的给负。

\subsection{一个完整的无约束优化演示}

设
\[
f(x,y)=x^2+xy+y^2-3x.
\]

梯度：
\[
\nabla f=
\begin{bmatrix}
2x+y-3\\x+2y
\end{bmatrix}.
\]
令梯度为零，解方程组：
\[
\begin{cases}
2x+y=3,\\
x+2y=0,
\end{cases}
\]
得 \((x,y)=(2,-1)\)。只有一个驻点。

Hessian：
\[
H=\begin{bmatrix}2&1\\1&2\end{bmatrix}.
\]
这个矩阵不含 \(x,y\)——在整个平面上处处相同。特征值为 \(\lambda=1,3\)，都是正数，所以 Hessian 处处正定。\((2,-1)\) 是严格局部极小。

更进一步，因为函数是正定二次型 \(x^2+xy+y^2=\frac12[(x+y)^2+x^2+y^2]\) 加线性项 \(-3x\)，当 \(\|(x,y)\|\to\infty\) 时函数值向正无穷远增长，所以 \((2,-1)\) 同时也是全局极小。

两点提醒：Hessian 的正定性只给局部结论，要升级为全局最优需要凸性、紧致可行域或者其他整体结构。另外，局部极小不等于``周围所有点都比它大''——严格的极小也只是在某个小邻域内最小，走出那个邻域可能有更低的地方。

\subsection{退化与边界：二阶检验无结论之后}

二阶检验对以下情况不给答案：
\begin{itemize}
\tightlist
\item Hessian 半正定或半负定（有零特征值但非零特征值同号）；
\item Hessian 为零矩阵（所有二阶导数为零）。
\end{itemize}

第一种情况的一个标志性例子是 \(f(x,y)=x^2-2xy+y^2=(x-y)^2\)。Hessian 在原点（其实在任意点）为
\[
\begin{bmatrix}2&-2\\-2&2\end{bmatrix},
\]
特征值 0 和 4，半正定。整条直线 \(x=y\) 都取相同的最小值零——这些点都是非严格极小。

第二种情况：\(f(x,y)=x^4+y^4\) 在原点 Hessian 为零矩阵，二阶检验完全沉默。但原函数非负且只在原点为零，仍有严格极小。此时四阶项主导了局部行为。

退化的排查方向：先看 Hessian 零空间里的方向，沿这些方向检查三阶、四阶项；如果所有方向都由二阶项覆盖了曲率判断，问题基本定位。

\subsection{Newton 步：把二阶模型最小化}

在 Hessian 可逆处，把 Taylor 展开截到二阶作为局部二次模型：
\[
q(\mathbf h)=f(\mathbf x)+\nabla f(\mathbf x)^{\trans}\mathbf h
+\frac12\mathbf h^{\trans}H_f(\mathbf x)\mathbf h.
\]
\(q(\mathbf h)\) 对 \(\mathbf h\) 的梯度为 \(\nabla f(\mathbf x)+H_f(\mathbf x)\mathbf h\)。令其为零，解得
\[
\mathbf h=-H_f(\mathbf x)^{-1}\nabla f(\mathbf x).
\]

更新规则 \(\mathbf x\leftarrow\mathbf x+\mathbf h\) 就是多元 Newton 法的核心。它在靠近非退化解时收敛极快（二次收敛），但离解太远或 Hessian 不定时，Newton 步可能走向极大甚至鞍点，而不是极小。在优化实践里，通常需要线搜索或信赖域来约束步长，或者对 Hessian 做正则化（加一个正的对角矩阵，使搜索方向折向负梯度）。

公式 \(\mathbf h=-H^{-1}\nabla f\) 来自局部二次模型——它不是全局安全结论，而是一个基于局部曲率信息的搜索建议。

本篇的二阶分析止于无约束驻点。当变量受到等式或不等式约束时，拉格朗日乘数法和 KKT 条件接手——那会改变``梯度为零''的涵义，在凸优化部分详细展开。
